\specialsection{Введение}


Объём образовательного видеоконтента на открытых платформах растёт быстрее, чем появляются инструменты для его оценки.
Пользователь, выбирая между десятками лекций по одной теме, вынужден ориентироваться на лайки,
просмотры и удержание, то есть метрики, разработанные для развлекательного контента и оптимизированные под клик, а не под передачу знаний.
Автор, в свою очередь, не получает содержательной обратной связи: что именно в его лекции не работает и для кого.

Ситуация воспроизводит логику «рынка лимонов» \cite{akerlof1970lemons}: обе стороны действуют в условиях информационной асимметрии.
Устойчивых сигналов, позволяющих сопоставить видео с образовательными целями конкретной аудитории, не существует.

Образовательное видео как объект оценки принципиально отличается от развлекательного: длинный таймлайн, высокая смысловая нагрузка, мультимодальность (речь, слайды, жесты),
неравномерность интерактивности и критическая зависимость от контекста целевой аудитории. Ни одна из классических платформенных метрик эти характеристики не учитывает.

Существующие инструменты решают смежные задачи: Gemini позволяет суммаризировать видео по ссылке и генерировать таймкоды, StudyFetch — формировать конспекты и флеш-карты.
Однако ни один из них не оценивает видео относительно конкретного пользовательского запроса и не даёт интерпретируемой аналитики для принятия решений.

В данной работе предлагается система автоматизированной оценки образовательного видеоконтента, формирующая аналитическую выдачу в зависимости от профиля пользователя:
его роли, уровня подготовки, погружённости в тему и цели просмотра. Система оценивает видео по двум осям — форма и содержание, и включает в себя техническое качество аудио и видео,
систему бинарных флагов восприятия на основе типологии швейцарского психолога Карла Густава Юнга, профиль подачи (академичность и инструментальность),
адаптивную нарративную суммаризацию с тематической сегментацией.
Все оценки формируются без эталонного референса — исключительно на основе анализа самого видео.

\textbf{Цель работы} — разработать инструмент автоматизированной оценки качества и релевантности образовательного видеоконтента,
формирующий интерпретируемую аналитику для обучающихся и авторов контента.

\vspace{20pt}

\textbf{Задачи:}
\begin{enumerate}
\item Провести обзор существующих подходов к оценке качества аудио-\nobreak\hspace{0pt}визуального контента и выявить их ограничения применительно к образовательному видео.
\item Разработать структуру каскадного анализа, охватывающего технические, поведенческие и содержательные аспекты.
\item Реализовать модули оценки: технического качества, флагов восприятия, профиля подачи и нарративной суммаризации.
\item Разработать пользовательский интерфейс, обеспечивающий интерпретируемое представление результатов анализа, 
в соответствии с выбранной ролью оценивающего: обучающийся, спикер, продакшн и заказчик.
\end{enumerate}

Программная реализация прототипа доступна в репозитории проекта~\footnote{\url{https://github.com/maxim-druzhinin/EduVisLab}}.
Репозиторий содержит код модулей пайплайна, зависимости и конфигурацию запуска. 